\cotefile{157-1}
\scanpage{114}
\mspage{120}
Par \unsure{contre}, on a

\uline{Prop. 8} Soit $f:X\to Y$ dans \unsure{Diag}. Conditions équivalentes

a) $f$ est une $\unsure{\mathbb D|D}$-fibration

b) \unsure{Pour} tout $p:Y''\unsure{\mathord{-}|\to}Y'$ sur $Y$, si $p$ est une $\mathbb D$-équiv. $1$-complète \unsure{(}déf. 1 p.~113), alors $q:X''-X'$ est une $\mathbb D$-équivalence.

% diagram in the left margin beside b)
\diagram{\begin{tikzcd}[ampersand replacement=\&,column sep=small,row sep=small]
X\arrow[d,"f"'] \& X'\arrow[l]\arrow[d,"f'"'] \& X''\arrow[l,"q"']\arrow[d,"f''"]\\
Y \& Y'\arrow[l] \& Y''\arrow[l,"p"]
\end{tikzcd}}

\uline{Cor.} Supposons que dans Diag, toute $\mathbb D$-équivalence soit $1$-complète. (NB Je ne connais que le dérivateur HOT pour satisfaire à cette condition, qui \unsure{paraît} quasiment \unsure{caractéristique} de HOT.) Alors $f:X-Y$ est une $\mathbb D$-\deleted{équiv.}fibration \unsure{ssi} pour tout $p:Y''-Y'$ sur $Y$, si $p$ est une $\mathbb D$-équiv., alors $q:X''-X'$ aussi. Dans ce cas ``exceptionnel'' parmi les dérivateurs, la notion de $\mathbb D$-fibration est stable par composition (ce qui est \unsure{tautologique|tout le gag} sur le critère que je viens de donner.)

On va renforcer la notion de $\mathbb D$-fibration, pour la rendre plus souple, et notamment stable par composition.

\uline{Proposition 9} Soit $f:X-Y$ dans Diag. Conditions équivalentes :

a) $\forall\,y\in Y,\quad X_y-X_{/y}$ est une $\mathbb D$-\uline{équiv. \unsure{cohomologique}} (déf. 1 p.~114)\unsure{''|,} et \unsure{de|en} $=$ pour $X_{/y}-X_{/y'}$ pour toute flèche $y\longrightarrow y'$ dans $Y$. Et ces conditions sur $f$ restent valables après tout chgt de base (\uline{NB} il suffit que la première de ces conditions le reste.)\unsure{.}
% archive sheet number: 114
\scanpage{115}
\mspage{121}
b) Pour tout $F\in\operatorname{Ob}\unsure{\mathbb D|D}(\unsure{x|X})$ \uline{loc.} \uline{constant}, $f_*F$ est loc. constant sur $Y$ et se calcule ``fibre par fibre'' -- et \unsure{itou|idem} après tout changement de base

c) Pour tout $Y''\underset{p}{\unsure{\mathord{-}|\to}}Y'$ sur $Y$, si $p$ est une $\mathbb D$-équiv. cohomologique, alors il en est de même de $q:X''\to X'$.

c$'$) Comme c, mais on suppose \deleted{que} comme hyp. sur $p$ que $Y''$ soit la cat. finale, et $Y'$ ait un objet final.

\uline{Dém} Dans l'ordre ($\pm$ \unsure{tautologique})

\[a\Longrightarrow b\Longrightarrow c\deleted{\unsure{s|'}}\Longrightarrow c'\Longrightarrow a\]

\uline{Définition 5}. Dans les conditions équivalentes de la prop. 9, on dit que $f$ est une $\mathbb D$-\uline{fibration} \uline{cohomologique}. Sous les conditions duales, on dit que $f$ est une $\mathbb D$-\uline{fibration homologique}. On dit que $f$ est une $\mathbb D$-\uline{fibration forte}, si c'est à la fois une $\mathbb D$-fibration cohomologique et une $\mathbb D$-fibration homologique.

Ces notions \deleted{\unsure{toutes deux}} renforcent la notion de $\mathbb D$-fibration. \inserted{Elles ont \unsure{l'avantage} de plus d'être \uline{stables par composition}, \unsure{ainsi que|car il suffit} \unsure{par|pour} chgt. de base.}

\deleted{\unsure{La} \illegible{2} à l'aide de ces deux \unsure{notions} le résultat \unsure{suivant} :}

\deleted{\uline{Prop. 10} Soit $f:X\to Y$ dans \unsure{Diag}. Conditions équivalentes :}

\deleted{a) $\forall\, y\in\operatorname{Ob}Y$, $X_y-X_{/y}$ est une $\mathbb D$-équiv. $1$-complète (resp. $2$ complète), et de $=$ pour $X_{/y}-X_{/y'}$ pour toute flèche $y-y'$ dans $Y$. Et ça reste valable après tout chgt de base (NB il suffit de l'exiger pour la première de ces conditions, concernant $X_y-X_{/y}$.)}

\deleted{\unsure{a$'$|b}) Condition duale de a, en termes des $X_y-{}_{y\backslash}X$, $\,{}_{y'\backslash}X\to{}_{y\backslash}X$, qui doivent être \unsure{tous|t.} des $\mathbb D$-équiv. $1$-complètes (resp. $2$-complètes), et ça reste vrai après tt chgt. de base.}

\deleted{\unsure{b)}}
% archive sheet number: 115
\scanpage{116}
\mspage{122}
\deleted{b) Pour tout $p:Y''\unsure{\mathord{-}|\to}Y'$ sur $Y$, si $p$ est une $\unsure{\mathbb D|D}$-équiv. $1$-compl. (resp. $2$-complète), il en est de même de $q:X''\to X'$.}

\deleted{b$_1^{\unsure{*|'}}$) Comme b$_1$), mais en supposant $Y''$ cat. \unsure{ponctuelle}, $Y'$ \deleted{\unsure{d'un}} un objet final}

\deleted{b$_2$) Comme b), mais en supposant $Y''$ cat. ponctuelle, $Y'$ un objet initial.}

\deleted{Je ne \unsure{connais} pas de \unsure{renforcement} co}

La notion de $\mathbb D$-fibration forte a été \deleted{\unsure{prise}} choisie pour \unsure{coiffer} les deux notions duales de \unsure{$\mathbb D$-fib.|$f$-fib.} coh. et hom. (à défaut d'une meilleure notion), tout en restant stable par composition. On pourrait exiger que pour tt

\[Y''\xrightarrow{p}Y'\]

sur $Y$, avec $Y',Y''$ \unsure{Hot-asphériques|tlt-asphériques}, $q:X''-X'$ soit une $\mathbb D$-équivalence $1$-complète (\unsure{voire} même, \deleted{\unsure{une}} $2$-complète), appeler ça les $\mathbb D$-\uline{fibrations} $1$-\uline{complètes} (resp. $2$-\uline{complètes}). \deleted{\unsure{Ça coiffe}} Renforce \deleted{\unsure{quel}} de façon assez élégante la notion de $\mathbb D$-fibration forte, mais a le grand tort de ne pas être stable par composition (que je sache !) ; sauf dans le cas particulier où on a des \unsure{Hom} internes \unsure{sympa} (cf bas de la page 113), auquel cas la notion de $\mathbb D$-équiv. coh. implique déjà celle de $\mathbb D$-équiv. $1$-complète,
% archive sheet number: 116
\scanpage{117}
\mspage{123}
donc aussi celle de $\unsure{\mathbb D|D}$-équiv. homologique. Donc dans ce cas

\diagram{\begin{tikzcd}[ampersand replacement=\&,column sep=large]
\mathbb D\text{-fibration coh.}\arrow[r,Leftrightarrow]\arrow[d,Rightarrow] \& \mathbb D\text{-fibration }1\text{-complète}\\
\mathbb D\text{-fibration hom.} \& 
\end{tikzcd}},

et il n'y a pas à chercher plus loin que la $\mathbb D$-fibration coh.

On pourrait songer, il est vrai, à définir \deleted{\unsure{même dans ce cas}} \inserted{(dans le cas général)} une notion plus forte, en exigeant que pour tout $p:Y''\unsure{\mathord{-}|\to}Y'$ \inserted{sur $Y$} qui soit une $\mathbb D$-équiv. $1$-complète, $q:X''-X'$ est aussi une $\mathbb D$-équiv. $1$-complète. Cette fois-ci, la notion est stable par composition, mais elle me paraît néanmoins inutilisable (en dehors du cas qu'on vient de dire, où elle coïncide avec la notion de $\mathbb D$-fibration coh., ou encore de $\mathbb D$-fibration forte), car n'étant pas en mesure de décrire les $\mathbb D$-équiv. $1$-complètes, je suis aussi incapable de \deleted{\unsure{en décri|en décr}} donner des exemples \inserted{non triviaux} de $\mathbb D$-fibrations ``strictement $1$-complètes''.

\figure{\begin{tikzpicture}\draw (0,0) -- (.19,-.24);\end{tikzpicture}}

Pour vérifier qu'\unsure{une|un} $f:X-Y$ dans Diag est une $\mathbb D$-fibration cohomologique, il faut
% archive sheet number: 117
\scanpage{118}
\mspage{124}
savoir vérifier que des flèches $\unsure{\phi|q}:Z'\unsure{\mathord{-}|\to}Z$ dans Diag sont des $\unsure{\mathbb D|D}$-équiv. cohomologiques, donc \deleted{\unsure{en fait que ce}} \inserted{que}

\[H^{\unsure{\cdot|'}}_{\mathbb D}\unsure{(|[}Z,F)\xrightarrow{\sim}H^{\cdot}_{\mathbb D}(Z',\phi^*(F))\]

est iso pour tout $\unsure{F|\mathbb F}\in\operatorname{Ob}\mathbb D(Z)$ qui est loc \unsure{constante} sur $Z$. Or je ne sais vérifier un énoncé de cette nature (en dehors du cas où je sais d'avance que $Z$ est complètement $\mathbb D$-asphérique, i.e. que les coeff. loc. constants sur $Z$ sont en fait constants) \uline{que} si $\unsure{q|\phi}$ est déjà une Hot-équivalence. C'est dire que je ne sais \deleted{décrire} \inserted{exhiber} des $\mathbb D$-fibrations cohomologiques, que quand ce sont déjà des HOT-fibrations cohomologiques. Je veux, dans la suite, donner des critères utilisables pour des HOT-fibrations. Mais on va axiomatiser \deleted{\unsure{ceci}} ces développements, en gardant un $\mathbb D$ général, sur lequel je vais faire une hypothèse \inserted{si} draconienne que pratiquement, je ne connais que le dérivateur HOT qui y satisfasse. Il s'agit de \unsure{ce|ces} \unsure{ci-dessous}

\uuline{Der 3} c$_1$) Toute $\mathbb D$-équivalence est une $\mathbb D$-équiv. cohomologique

c$_2$) Toute $\mathbb D$-équivalence est une $\mathbb D$-équiv. \unsure{homologique}
% archive sheet number: 118
\scanpage{119}
\mspage{125}
$c_{\deleted{\unsure{1|2}}}'$) Toute $\unsure{\mathbb D|D}$-équivalence est \deleted{\unsure{une}} $1$-complète

$c''$) Toute $\mathbb D$-équivalence est $2$-complète.

\uline{NB} On a

\diagram{\begin{tikzcd}[ampersand replacement=\&,row sep=small,column sep=small]
\& c''\arrow[d,Rightarrow] \&\\
\& c'\arrow[dl,Rightarrow]\arrow[dr,Rightarrow] \&\\
c_1 \& \& c_2
\end{tikzcd}}

\unsure{tautologiquement}. Je ne connais que HOT pour satisfaire \deleted{\unsure{les}} \inserted{\unsure{à l'une|l'une} des deux} plus faibles conditions $c_1,c_2$. Et en fait, sous réserve d'écrire la théorie du dérivateur HOT, il doit être vrai qu'elle satisfait même la forme la plus forte $c''$. (Et les HOT-équiv doivent être simplement les Hot-équivalences.) Cf le ``lemme'' p.~\unsure{104|109}.)

La proposition suivante est une tautologie

\uline{Proposition 10} Supposons que $\mathbb D$ satisfasse la condition Der 3 $c_1$ (p.~124), p. ex. $\mathbb D=\mathrm{HOT}$. Alors \unsure{la|les} notion de $\mathbb D$-fibration coïncide avec celle de $\mathbb D$-fibration cohomologique. Elle est donc stable par composition. \inserted{Pour que $f$\inserted{$:X\unsure{\mathord{-}|\to}Y$} soit une $\mathbb D$-fibration il \unsure{f. et s.} que pour $Y''\xrightarrow{p}Y'$ sur $Y$, si $p$ est une $\mathbb D$-équiv., il en est de même pour $q\unsure{=|:}X''\to X'$.}

\uline{Corollaire 1} Si $f:X\to Y$ dans Diag est $\mathbb D$-propre \deleted{(resp. $\mathbb D$-lisse)}, alors $f$ est une $\mathbb D$-fibration cohomologique \unsure{sss|ssi} \deleted{pour} toute flèche $y'-y$ dans $Y$, $X_{/y'}-X_{/y}$ est une $\mathbb D$-équiv. (i.e. sss $f$ est une $\mathbb D$-fibration \inserted{coh.} faible).

\unsure{Appliquons} le critère c) de prop. \unsure{6|5}.
% archive sheet number: 119
\scanpage{120}
\mspage{126}
\uline{Cor. 2} Supposons que $\unsure{\mathbb D|D}$ satisfasse Der 3 $(c_1,c_2)$ (\unsure{par|p.} ex. $\mathbb D=\mathrm{HOT}$). Alors \deleted{dans $\mathbb D$} pour une flèche $f:X\unsure{\mathord{-}|\to}Y$ dans Diag, il revient au $=$ que $f$ soit une \deleted{fib} $\mathbb D$-fibration, une $\mathbb D$-fibration cohomologique, ou une $\mathbb D$-fibration homologique. Si $f$ est propre (resp. lisse), \unsure{il} revient au $=$ de dire que $f$ est une $\mathbb D$-fibration coh. (resp hom.) faible.

\figure{\begin{tikzpicture}\draw (0,0) -- (.2,-.23);\end{tikzpicture}}

\deleted{\unsure{Je}}

On voit donc que dans Cat, on doit avoir une seule notion \unsure{vraisemblable|raisonnable} de HOT-fibration, qui \unsure{coïncidera} avec celle de HOT-fibration cohomologique ou homologique. Mais pour que cette notion soit dès à présent disponible, il faudrait que le dérivateur HOT soit déjà construit, sur le domaine $\mathrm{Diag}=\mathrm{Cat}$ tout entier. Or ce n'est pas chose faite encore (et il n'a \deleted{\unsure{donc}} pas pu être question de prouver les propriétés de HOT dont il vient d'être question, disons la plus forte Der 3 $c''$).

Néanmoins, je voudrais développer de façon directe (``élémentaire'') la notion de Hot-fibration, sans me référer à la théorie du dérivateur HOT, et aux développements théoriques précédents,
% archive sheet number: 120
\scanpage{122}
\diagram{\begin{tikzcd}[ampersand replacement=\&,column sep=small,row sep=huge,execute at end picture={\draw[rounded corners=3pt] ([xshift=-5pt,yshift=9pt]\tikzcdmatrixname-1-1.north west) rectangle ([xshift=11pt,yshift=-6pt]\tikzcdmatrixname-2-2.south east);
\node[anchor=south] (sharedlabel) at ([xshift=3pt,yshift=10pt]\tikzcdmatrixname-1-2.north east) {$\scriptstyle\unsure{\mathrm{bigr}^{\mathrm o}|\mathrm{topos}|\mathrm{bij}^{\mathrm n}}$};
\draw[->] (\tikzcdmatrixname-1-2.north east) to[bend right=12] (sharedlabel.south west);
\draw[->] (\tikzcdmatrixname-2-2.east) to[bend right=34] (sharedlabel.south);}]
\unsure{\mathscr H|\mathcal H}(A)\arrow[r,"{\deleted{\unsure{\varphi_A}}\unsure{L_A|\varphi_A}}"]\arrow[d,"f_!"']\arrow[rrr,bend left=28,"{\int_A}"]
\& \mathscr H_{\unsure{\mathrm{top}}}(A^\wedge)\arrow[d,"{(f_{\wedge}^{\ast})^{\ast}}"]\arrow[r,"{\alpha_{A^\wedge}}"]
\& \underline{\operatorname{Hom}}(A^\wedge,H)\arrow[d]\arrow[r,"{\varepsilon_{A^\wedge}}","{\psi\unsure{\mapsto|\vdash}\psi(e_{A^\wedge})}"']
\& H\arrow[d,equal]\\
\mathscr H(B)\arrow[r,"{\unsure{\varphi_B|\psi_B}}"] \& \mathscr H_{\mathrm{top}}(B^\wedge)\arrow[r]
\& \underline{\operatorname{Hom}}(B^\wedge,H)\arrow[r,no head,"{\varepsilon_{B^\wedge}}"] \& H
\end{tikzcd}}

\[\unsure{\mathscr H(\,)|\mathscr H(1}\]

\diagram{\begin{tikzcd}[ampersand replacement=\&,column sep=small,row sep=small]
\deleted{\unsure{X}}\arrow[d,no head,"\deleted{\unsure{f}}"'] \&\\
A\arrow[d] \& X\arrow[l]\\
B \& Y\arrow[l]
\end{tikzcd}}

\[\operatorname{Hom}_B(f_!(X),Y)\simeq\operatorname{Hom}_A(X,f^*(Y))\]

[$Y$ cofibré sur $B$, de sorte que $f^*(Y)$ est bien donné par $A\times_B Y$]

valable également en passant à $\unsure{\mathrm{Hot}(B)|\mathrm{Hot}-(B)|\mathrm{Hot}(A)}$

\[\boxed{\operatorname{Hom}_{\mathscr H(B)}(f_!(X),Y)\simeq\operatorname{Hom}_{\mathscr H(A)}(X,f^*(Y))}\]
% archive sheet number: 122
\scanpage{124}
\diagram{\begin{tikzcd}[ampersand replacement=\&,column sep=small,row sep=small,cells={font=\scriptsize}]
\&\&\text{fibré parfait}\arrow[dl]\arrow[d]\arrow[dr]\&\&\\
\&\text{propre.}\arrow[d]\&\text{fibrant}\arrow[dl]\arrow[dr]\&\text{lisse}\arrow[d]\&\\
Y\to\mathrm{Hot}_{\unsure{\mathbb D|D}}\&\text{préfibr. g.}\&\&\text{préfibr. dr.}\&Y^{\unsure{\mathrm o|0}}\to\mathrm{Hot}_{\mathbb D}
\end{tikzcd}}

\diagram{\begin{tikzcd}[ampersand replacement=\&,row sep=small,column sep=small]
\&\mathbb D\text{-équiv.}\&\\
\mathbb D\text{-coasph.}\arrow[ur]\&\&\mathbb D\text{-asph.}\arrow[ul]
\end{tikzcd}}
\noindent\rule{.68\linewidth}{.45pt}

\begin{center}Le Ciel\end{center}

\diagram{\begin{tikzcd}[ampersand replacement=\&,row sep=small,column sep=2pt,cells={font=\scriptsize},execute at end picture={\draw[gray,line width=.8pt] ([xshift=-2pt]\tikzcdmatrixname-3-1.west) to[bend left=14] (\tikzcdmatrixname-1-3.center) to[bend left=13] ([xshift=2pt]\tikzcdmatrixname-3-5.east) -- ([xshift=3pt,yshift=-5pt]\tikzcdmatrixname-4-5.south east) -- (\tikzcdmatrixname-4-5.south west) -- (\tikzcdmatrixname-4-3.north east) -- (\tikzcdmatrixname-4-3.north west) -- (\tikzcdmatrixname-4-1.south east) -- ([xshift=-2pt]\tikzcdmatrixname-4-1.south west) -- cycle;
\draw[gray,->] ([xshift=17pt,yshift=9pt]\tikzcdmatrixname-4-5.east) -- (\tikzcdmatrixname-6-3.north east);}]
\text{cohomologie}\&\&\text{\deleted{Le Ciel}}\&\&\text{homologie}\\
\&\&\text{Cat-bifibrant}\arrow[d,Rightarrow]\arrow[dll,Rightarrow]\arrow[drr,Rightarrow]\&\&\\
\substack{\text{\deleted{cohomologique}}\\\text{Cat-cofibrant}}\arrow[d,Rightarrow]\&\&\mathbb D\text{-parfait}\arrow[dll,Rightarrow]\arrow[d,Rightarrow]\arrow[drr,Rightarrow]\&\&\substack{\text{\deleted{homologique}}\\\text{Cat-fibrant}}\arrow[d,Rightarrow]\\
\mathbb D\text{-propre}\arrow[dd,Rightarrow]\&\&\mathbb D\text{-}\unsure{\text{fibré fort}|\text{fibrat. forte}}\arrow[dl,Rightarrow]\arrow[ddll,Rightarrow]\arrow[dd,Rightarrow]\arrow[ddrr,Rightarrow]\arrow[dr,Rightarrow]\&\&\mathbb D\text{-lisse}\arrow[dd,Rightarrow]\\
\&F'\arrow[dl,Rightarrow]\&\&F''\arrow[dr,Rightarrow]\&\\
\mathbb D\text{-préfibrant coh.}\arrow[d,Rightarrow]\&\&\mathbb D\text{-fibré}\deleted{\unsure{\text{fort}}}\arrow[dll,Rightarrow]\arrow[drr,Rightarrow]\&\&\mathbb D\text{-préfibrant hom.}\arrow[d,Rightarrow]\\
\mathbb D\text{-préfibrant à g}\&\&\unsure{\uparrow|1}\&\&\mathbb D\text{-préfibr. à droite}
\end{tikzcd}}

% slanted
\margin{\unsure{$=$ notion} dans HOT}
% left brace beside the last two nodes on the left; slanted
\margin{$\bigl\{$ \unsure{$=$ notion} dans HOT}
% right brace beside the last two nodes on the right
\margin{$\bigl\}$ \unsure{$=$ notion} dans HOT}

\deleted{\unsure{S}}

\[\begin{aligned}F'&=\mathbb D\text{-fibrations cohomologiques}& F''&=\mathbb D\text{-fibrations homologiques}\end{aligned}\]

Si $f:X\unsure{\mathord{-}|\to}Y$ est \uline{pré}fibrant coh. et hom., \unsure{alors|donc} faut il \uline{fibrant} coh. et hom ?

\[\mathbb D\text{-équiv. coh}=\mathbb D\text{-équiv hom}\]

\margin{OK si\\ Est-ce tjrs le cas ?\\ J'en doute.}

Est-ce équivalent \deleted{pour} $\mathrm{Der}(k\text{-}\mathrm{Mod})$, $k$ un anneau ?
% archive sheet number: 124
\scanpage{125}
\mspage{103}
% letterhead
Alain LAMBIN DOSTROMON\\
Lo Marchadiol\\
19500 TURENNE\\
(France)\\
Tel : 55.85.90.23\\
19 1 00588 3\\

\handwritten{10 . 4 \unsure{86|76}}\\

\handwritten{Pour mieux vous défendre, voici un courrier adressé par moi même à la \unsure{Sécu|Sécu.} de la \unsure{Correze|Corrèze} à propos \unsure{de|du|des} \unsure{Sain|Sabin|Bain|soin} - ceci est valable pour \uuline{tous} les vaccins -}

\handwritten{Bon courage -}\\
\handwritten{de tout cœur avec}\\
\handwritten{vous.}\\

\handwritten{\unsure{A|AL}}\\

% printed form
En cas de courrier, joindre une enveloppe timbrée à votre adresse, merci.

% printed form
Veuillez garder la présente ordonnance, je me vois dans l'obligation de ne plus faire de duplicata

% printed form
Gardez l'original, \unsure{expédiez|expediez} la copie certifiée à votre Caisse

% printed form
Membre d'un centre de gestion \unsure{agréé} par l'Administration \unsure{fiscale}, acceptant les \unsure{règlements|réglements} par chèques \unsure{libellés} à mon nom
% archive sheet number: 125
\scanpage{126}
Variantes pour $X\to Y$ $\unsure{\mathbb D|D}$-propre, \unsure{i.e.} les

\[i_y:X_y\unsure{\mathord{-}|\to}X_{/y}\]

$\mathbb D$-coasphérique. On peut envisager

% slanted
\margin{$\mathbb D$-\unsure{propriété|propreté} faible\\ $\mathbb D$-propriété \deleted{faible}\\ \fbox{\shortstack{\uline{faiblement} $\mathbb D$-\uline{propre}\\ i.e. \unsure{quasi} $\mathbb D$-propre}}\\ $\Downarrow$ Pour \unsure{HOT|Hot}, les \unsure{conditions} a) b) c) d) coïncident.}

\[\unsure{e}\xrightarrow{\text{coasph.}}Y_{/y}\]

\deleted{\unsure{$f^*f_*$}}

\uline{$f_*f^*M$ \inserted{\uline{loc. const.}} se calcule fibre par fibre.} stable par chg. de base \uline{lisse}

[\unsure{0|a}) les $i_y$ $\mathbb D$-coasphériques ]

\textcircled{\scriptsize\unsure{a|a$'$|b}} les $i_y$ des $\mathbb D$-équiv -- \deleted{\unsure{mais}} et c'est vrai après tt chgt de base

\margin{$\Uparrow$ Suffit pour\\ $Y-\mathrm{Hot}_{\mathbb D}$}

\[\boxed{\text{$f$ \unsure{préfibrant} à $\mathbb D$-\unsure{dr}}}\]

$\Longleftrightarrow$ ? les $f_*(F)$ ($F$ \unsure{cst}) se calculent fibre par fibre

Si

\diagram{\begin{tikzcd}[ampersand replacement=\&,row sep=small,column sep=small]
X''\arrow[r,no head,"q"]\arrow[d,no head] \& \unsure{X'|X}\arrow[d,"f'"'] \& X\arrow[l,<->]\arrow[d]\\
Y''\arrow[r,no head,"p"']\& Y'\arrow[r]\& Y
\end{tikzcd}}

avec $p$ \uline{coasph.}, alors $q$ est $\mathbb D$-équiv. ??

b) les $i_y$ des $\mathbb D$-équiv. coh., et stables par chgt. de base

\deleted{\unsure{$\Longleftrightarrow$}}

$\Longleftrightarrow$ les $f_*(F)$ ($F$ loc. ct) se calculent fibre par fibre, et idem après chgt de base

\margin{$\Uparrow$ suffit pour\\ $Y-\mathrm{Hot}_{\mathbb D,\mathrm{coh.}}$}

\[\boxed{\mathbb D\text{-préfibration cohomologique}}\]

$\Longleftrightarrow$ si $p$ comme dessus, $q$ est \inserted{$\mathbb D$-}équiv coh.

\diagram{\begin{tikzcd}[ampersand replacement=\&,column sep=small,row sep=small]
X'\arrow[r,no head,"q"]\arrow[d,no head,"f'"']\& X\arrow[d,no head,"f"]\\
Y'\arrow[r,"p","{\text{coasph.}}"']\&Y
\end{tikzcd}}

$\deleted{\unsure{\mathrm{eq}}}$

\[H^{\unsure{\cdot|'}}_{\mathbb D}(X,\unsure{a_X|o_X})\simeq H^{\cdot}_{\mathbb D}(Y,\underbrace{f_*(a_X)}_{\deleted{\unsure{\text{loc. const.}}}})\simeq H^{\cdot}_{\mathbb D}(Y',\underbrace{p^*f_*\unsure{(a_X)}}_{f'_*\underbrace{q^*(a_X)}_{a_{X'}}})\]

$F-$ \qquad $\deleted{\unsure{e\to Y}}$ \deleted{\unsure{objet}}

$\deleted{\unsure{\simeq H_{\mathbb D}}}$ \qquad $\unsure{y\backslash e}$

\margin{$\Uparrow$}

\unsure{c}) les $i_y$ des $\mathbb D$-équiv. $1$-complètes et \deleted{\unsure{le reste}} après chgt de base

\[Y\to\mathrm{Hot}_{\mathbb D,1}\quad\deleted{\unsure{\mathbb D\text{-fibré cohomologique}}}\]

$\Longrightarrow$ \tikz[baseline=(qc.base)]\node[draw,ellipse] (qc) {\shortstack{$q$ est $\mathbb D$-\unsure{équ.} $1$-\unsure{cpt.} ?\\???}};

\uline{NB} les \deleted{\unsure{$f$}} coasphériques sont des $\mathbb D$-équiv. \uline{coh} ??? oui, et $=$ des $\mathbb D$-équiv. $1$-\uline{complètes} (et $=$ \deleted{\unsure{de 2}} \uline{complètes} ?)

\margin{$\Uparrow$}
\textcircled{\scriptsize d}) idem avec $2$-\uline{complète}

\[Y\to \tikz[baseline=(hc.base)]\node[draw,ellipse,inner sep=2pt] (hc) {$\mathrm{Hot}_{\mathbb D,2}$};\qquad\boxed{f\ \mathbb D\text{-}\deleted{\text{\unsure{faiblement}}}\text{quasi-propre}}\]

\deleted{\unsure{$\times$}} ? \qquad $\deleted{\unsure{F,G}}$ \qquad $\deleted{\operatorname{Hom}(F,G)}$

\margin{$\Uparrow$}
\textcircled{\scriptsize\unsure{e}} \unsure{idem} avec coasphérique $=$ propre \qquad $Y\to\mathrm{Hot}_{\mathbb D\,2}$

$\Longleftrightarrow$ \deleted{\unsure{princ.}} \unsure{!} Si $p$ coasph, $q$ aussi

Si $f$ est faiblement $\mathbb D$-propre et faiblement $\mathbb D$-lisse, c'est une $\mathbb D$-fibration.
% archive sheet number: 126
\scanpage{127}
\mspage{105}
% letterhead
\textsc{Docteur Guy} DESMICHELLE\\
PROFESSEUR AU CENTRE HOMÉOPATHIQUE\\
DE FRANCE\\
CHARGÉ DE COURS AUX FACULTÉS\\
DE MÉDECINE ET PHARMACIE\\
\rule{1cm}{.4pt}\\
68, \textsc{Bd Saint\unsure{.|-}Michel}, 75006 PARIS\\
43.26.99.75\\
\rule{1cm}{.4pt}\\
SUR RENDEZ\unsure{.|-}VOUS\\
\rule{1cm}{.4pt}\\

Paris, le 10 Mars \unsure{I|1}986\\

Monsieur Jean COULARDEAU\\
La Ribe - Les\unsure{.} Vastres\\
43430 FAY SUR LIGNON\\

Cher Monsieur,\\

Merci de m'avoir tenu au courant régulièrement de vos difficultés et bravo pour la sentence du Tribunal du Puy.

Les mentalités changent, on peut espérer qu'un jour la législation évoluera et que les administrations évolueront à leur tour.

Comme je crois vous l'avoir signalé il y a quelques années, je suis un ami d'enfance du Docteur Michel CHAZOT de MONTFAUCON et donc un ami de la famille CHAZOT et VEY de FAY.

Je vous prie de croire, Cher Monsieur, en mes sentiments\\
distingués et les meilleurs.\\

\handwritten{\illegible{1}\quad\tikz[x=1cm,y=1cm]{\draw (0,.27) .. controls (.04,-.02) and (.35,-.10) .. (.61,0) .. controls (.47,.15) and (.12,.29) .. (0,.27); \draw (.7,.18) .. controls (.81,-.01) and (.97,.02) .. (1.08,.05) .. controls (1.02,.18) and (1.11,-.02) .. (1.5,.04) .. controls (1.53,-.08) and (1.87,-.07) .. (2,.01) .. controls (1.74,.22) and (1.68,.11) .. (2.15,.02) .. controls (2.44,-.04) and (2.53,-.07) .. (2.69,-.05) .. controls (2.36,.19) and (2.28,.15) .. (2.65,.02) .. controls (3.01,-.11) and (3.83,-.29) .. (4.15,-.18) -- (4.17,-.69);}}\\
% archive sheet number: 127
\scanpage{128}
\[X\overset{f}{\unsure{\mathord{-}|\to}}Y\qquad f\text{ coasph}\]

$F$ loc. const. alors $F$ \unsure{sont} \deleted{\unsure{ctes}} sur les $X_{/y}$ par Der 3 a$'$

$f^*f_*F\to F$ iso ? c'est \unsure{local} sur $Y$, et OK.

\diagram{\begin{tikzcd}[ampersand replacement=\&,row sep=small,column sep=small,cells={font=\scriptsize},execute at end picture={
\draw[->] (\tikzcdmatrixname-3-2.south west) -- (\tikzcdmatrixname-4-1.east);
\draw[->] (\tikzcdmatrixname-3-2.south east) -- (\tikzcdmatrixname-4-3.west);
\draw[thin] (\tikzcdmatrixname-4-1.east) -- ([yshift=-3pt]\tikzcdmatrixname-3-2.south west) -- ([xshift=-3pt]\tikzcdmatrixname-4-1.east) -- ([xshift=5pt]\tikzcdmatrixname-3-2.south west) -- (\tikzcdmatrixname-3-2.center |- \tikzcdmatrixname-4-1.east) -- (\tikzcdmatrixname-4-3.west);
\draw[thin] (\tikzcdmatrixname-3-2.south) -- ([xshift=2pt,yshift=-5pt]\tikzcdmatrixname-4-3.west) -- (\tikzcdmatrixname-4-1.east) -- (\tikzcdmatrixname-5-2.center);}]
\unsure{\mathbb D|D}\text{-}\deleted{\text{co}}\text{asphérique}\arrow[dr,Rightarrow]\&\&\mathbb D\text{-coasphérique}\arrow[dl,Rightarrow]\\
\&\mathbb D\text{-équiv. }2\text{-compl.}\arrow[d,Rightarrow]\&\\
\&\mathbb D\text{-équiv. }1\text{-complète}\arrow[dl,Rightarrow]\arrow[dr,Rightarrow]\arrow[dd,Rightarrow,"{\deleted{\unsure{\mid}}}" description]\&\\
\substack{\mathbb D\text{-équiv.}\\\deleted{\mathbb D\text{-équiv. coh.}}\\\text{coh.}\\\text{compl.}}\arrow[ddr,Rightarrow]\&\&\substack{\deleted{\mathbb D\text{-}\unsure{\text{équi}}}\ \mathbb D\text{-équiv. hom}\\\text{complète}}\arrow[ddl,Rightarrow]\\
\&\deleted{\mathbb D\text{-équiv.}}\&\\
\&\mathbb D\text{-équiv.}\&
\end{tikzcd}}

\margin{\uline{On suppose Der 3 a$'$}}
% archive sheet number: 128
\scanpage{129}
\mspage{104}
% letterhead
docteur alain lambin-dostromon\\
lo marchadiol\\
turenne 19500 meyssac\\
tél. : (55) 85-90-23\\
sur rendez-vous\\

\margin{\unsure{XXX}}

Octobre 1985\\
à\qquad Mr SOLE\\
Directeur Adjoint CPAM\\

% slanted
\margin{COPIE}

Monsieur le Directeur,\\

Je reçois bien votre courrier du 2 septembre dernier relatif à la ``campagne de vaccination contre la grippe en faveur des personnes âgées''.

Je ne puis me soustraire au devoir que m'impose le code de déontologie, qui me fait obligation de ne pas garder secret ce que je sais.

Vous écrivez que ``cette expérience (de vaccination sur 3 ans) a permis une amélioration de la couverture vaccinale des personnes âgées''. Ceci est une évidence, mais ne signifie pas l'amélioration de l'état de santé de la population en question.

Je me dois de vous informer que tout vaccin allopathique a pour but d'augmenter les défenses d'un organisme déficient. D'où augmentation des gamma globulines ; \unsure{mais}... celles ci restent élevées devant l'agression vaccinale, et l'on crée ainsi un terrain tuberculinique, au sens de l'Homéopathie et de la Médecine Quantique.

Ceci va bientôt s'agraver d'une baisse des sérines, correspondant à une sycose. L'association des deux phénomènes crée une tuberculinosycose, ou cancérinisme, qui fait le lit.

- des cancers (vieillissement cellulaire et défaut de transfert des informations).

- des scléroses d'organes (insuffisance rénale, cardiaque, hépatique, diabète, artérite, artériosclérose, etc...) pour 85 \% des cas et aux

- scléroses nerveuses (débilité mentale, schizophrénie, sclérose en plaque, maladie de charcot, sclérose latérale, \unsure{amyotrophique}, etc...) pour 15 \% des cas.

Il existe d'autres évolutions inverses, sous vaccination, chez les sujets plus jeunes, avec hypo gamma globulinémie et hyper-sérine.

Ces conclusions, issues de nos travaux au sein de la Société Internationale de Médecine Quantique , et de l'utilisation de bilans protéiques informatisés, avec un recul d'au moins 10 ans pour nous même.

Je dois avouer que mon attention a été attirée sur les pathologies iatrogènes dues aux vaccins par le fait que je suis moi même une victime du vaccin antipolio injectable, réputé totalement atoxique... par les laboratoires qui le vendent !!! Et qui sont ceux que vous citez.

% letterhead
docteur alain lambin-dostromon\\
lo marchadiol\\
turenne 19500 meyssac\\
tél. : (55) 85-90-23\\
sur rendez-vous\\

Je ne critique pas leurs recherches, mais leurs méthodes commerciales.

Je me permets de vous citer mon expérience au Cabinet médical :

La protection par les vaccins que vous citez est assez correcte contre la grippe sauf :

- que les suites graves citées plus haut n'ont pas été mises en rapport avec la vaccination, dans l'esprit des allopathes. Peut être parce que ceux ci ne voient plus les maladies graves dont je parle.

- que les grippes graves déclanchées par le vaccin ne sont pas négligeables ni en nombre, ni quand aux conséquences vitales parfois, je suis surpris, aussi, que vous ne connaissiez pas les protections homéopathiques par Influenzinum ou Oscillococcinum selon l'âge du sujet.

Infiniment plus efficace et sans danger du type de ceux cités plus haut. Mon expérience de dix ans le prouve aisément, et se vérifie par tout médecin de bonne volonté.

Produits moins chers, pas dangereux, faciles d'emploi, baissans la consommation médicale. Il me semblerait plus logique de les intégrer à la liste des produits que vous citez.

Ainsi les personnes âgées auraient le choix, de leur thérapeutique, conformément aux droits de l'Homme, dont on parle d'autant plus qu'on les baffoue perfidement.

Pourriez vous avoir l'amabilité de me répondre. Ces courriers pouvant éventuellement être publiés.

Vous comprendrez aussi que, médicalement parlant, je ne puisse adhérer à une campagne vaccinale qui risque de créer encore plus de malades âgées, et ruiner un peu plus mes finances qui les alimentent, et le soignent.

Dans l'attente de votre réponse, je vous prie d'agréer, Monsieur\\
le Directeur, mes salutations les plus sincères.\\

\handwritten{\unsure{A|AL}}\\

PS : à votre disposition pour vous éclairer mieux sur la question.
% archive sheet number: 129
\scanpage{130}
Alors
\[
\unsure{\mathbb D|D}\text{-\unsure{Equivalence|Équivalence} coh.}
\underset{\substack{\vert\\\text{\unsure{astuce|existence}}\\\text{des }R\mathrm{Hom}}}{\Longleftrightarrow}
\mathbb D\text{-équiv. 1-complète}\Longrightarrow\mathbb D\text{-équiv. hom.}
\]

Mais a-t-on $\mathbb D\text{-équiv. hom}\overset{?}{\Longrightarrow}\mathbb D\text{-équiv. coh.}$ ?

Thomason a l'air de les \unsure{jeter|ficher} dans le $=$ sac ! lui poser la question !
% archive sheet number: 130
\scanpage{131}
Brives-Charensac, le 08/04/1986\\
Monsieur le Préfet,\\

Avec tous ceux qui oeuvrent pour une prise en charge personnelle de la santé, que cela soit celle des humains ou celle des animaux, les deux étant intimement liées (il est bien évident que les lois veillant à la protection du cheptel ont en vue la santé publique), les homœopathes se sentent fortement concernés par le procès qui est fait à Monsieur COULARDEAU pour refus de vacciner des animaux traités par homœopathie :

- Ce procès revient à nier purement et simplement, de la part des pouvoirs publics, l'homœopathie, sa valeur préventive, donc sa valeur tout \unsure{\uline{court}|court}, puis qu'elle se donne essentiellement comme une médecine de terrain.

- Il choque la conscience de tous ceux qui ont fait l'expérience du retour à une immunité naturelle solide après traitement homœopathique, et veulent en faire bénéficier leurs enfants.

- Il insulte les médecins qui se donnent le mal, pour le bien de leur prochain, d'acquérir cette difficile technique.

- Il choque tous les amis de la liberté de pensée, renouvelant une sorte d'inquisition qui cherche à obtenir de force l'adhésion de la conscience à des dogmes indémontrables.

Des événements récents publiés par la presse confirment que l'on ne touche pas impunément à l'appareil immunitaire d'un organisme vivant, et si les laboratoires scientifiques ne peuvent pas prouver (ou ne veulent pas) qu'une agression vaccinale puisse faire basculer un organisme en difficulté transitoire, ils ne peuvent pas non plus nous démontrer le contraire.

Doit-on, en république française, supporter de voir Mr COULARDEAU payer une amende pour refus d'appartenir à la religion d'état ?

Si la loi impose à chaque éleveur de maintenir les animaux en bonne santé, il semble prouvé que Mr COULARDEAU a satisfait à la loi plus que nombre d'éleveurs qui vaccinent leurs bêtes. Il a donc satisfait à l'esprit de la loi plus que ceux qui satisfont à la lettre.

\hfill .../...

% First copy of the closing, at left when the signature sheets are right side up.
En attendant que soit réformée la lettre au nom de l'esprit qui l'a inspirée, nous vous prions, Monsieur le Préfet, de bien vouloir user de votre autorité pour que Mr COULARDEAU soit jugé en fonction de cet esprit.

Dans cette confiance, nous vous prions de croire, Monsieur le Préfet, à l'expression de notre haute considération.\\

\handwritten{\unsure{A.|J.} Fayeton}\\

\handwritten{D\textsuperscript{r} Alain \unsure{Reynaud|Reynauld} \qquad \unsure{R}\\
(Montpellier)}\\

\handwritten{D \unsure{Diendomné|Dieudonné} \qquad \illegible{1}\\
le PUY.}\\

\handwritten{D\textsuperscript{r} Denis \unsure{GRANGEOUT|GRANGEON} \unsure{13} cours Victor HUGO 43000 LE PUY \qquad \unsure{Grangeon}}\\

\handwritten{Dr Damien DU PERRON \unsure{Ch|ch} \unsure{fr|St} Jean. St \unsure{Todard|Jodard}. \unsure{42590|42390} \unsure{NEULISE} \qquad \unsure{Du Perron}}\\

\handwritten{D\textsuperscript{r} \unsure{Naouri|Naoury} 60 Rue du P\textsuperscript{r} E Herriot Lyon \unsure{69002} \qquad \illegible{1}}\\

\handwritten{Dr G. \unsure{LOUTAN|COUTAN} 7. \unsure{Conseil-Général} 1205 GENÈVE \qquad \illegible{1}}\\

\handwritten{Dr P. \unsure{Harzé|Harzer} ch Ste \unsure{d'Hudard|d'Hudson} 42550 \unsure{Nervisé|Nervieux} \qquad \illegible{1}}\\

\handwritten{Dr. F. Vidal 34 rue de l'Horloge. 63200 Riom. \qquad \unsure{F V}}\\

\handwritten{Dr \unsure{Ch|de} Louise \unsure{Antony|Autton}, 26, rue de la \unsure{Griselle|Gabelle} 1228 \unsure{Genève|Genève-} Genève}\\

% Second copy of the closing, with a different set of signatures.
\mspage{102}
En attendant que soit réformée la lettre au nom de l'esprit qui l'a inspirée, nous vous prions, Monsieur le Préfet, de bien vouloir user de votre autorité pour que Mr COULARDEAU soit jugé en fonction de cet esprit.

Dans cette confiance, nous vous prions de croire, Monsieur le Préfet, à l'expression de notre haute considération.\\

\handwritten{D\textsuperscript{r} FAYETON. BRIVES CHARENSAC. 43700. \qquad \unsure{A.|J.} Fayeton.}\\

\handwritten{D\textsuperscript{r} Marc \unsure{DEGRAIX|DEGRAUX} de LAVISON \qquad \unsure{Alfred|Dfrid|David}\\
11 rue Luizet \qquad \unsure{d'Hillery|d'Hilley}\\
69130 ECULLY.\\
\unsure{13660} Orgon}\\

\handwritten{D\textsuperscript{r} LAURENT.A. \qquad \unsure{A Laurent}\\
13 Av. des 4 Chemins\\
92330 Sceaux}\\

\handwritten{DR \unsure{PETITJEAN}. Lyon \qquad \unsure{Petitjean}}\\

\handwritten{\unsure{SAUSSUS|SAUSSIS} \quad SEYSSINS 38170 \qquad \unsure{Saussus}}\\

\handwritten{Dr VIGNERON J.L CONDRIEU \qquad \unsure{Vigneron}}\\

\handwritten{D\textsuperscript{r} \unsure{Annie|Anne} \unsure{Jacquin|Jacquet} \unsure{Barre} \qquad \unsure{A J Barre}\\
Les \unsure{Sabines|Sabins} 69130 Ecully}\\

\handwritten{\unsure{Didier} \unsure{Cottencin|Cottereau} \qquad \illegible{1}\\
57 \unsure{bd|fbg} St Michel - 75005 Paris}\\

\handwritten{D\textsuperscript{r} \uline{\unsure{M. de K.|M. de K. S}}\\
32 rue St Joseph 31400\\
Toulouse}\\

\handwritten{D\textsuperscript{r} \unsure{MARZIALE|MARZIAKE} \qquad \illegible{1}\\
365 rue de Vaugirard\\
75015 Paris}\\

\handwritten{D\textsuperscript{r} \unsure{BARRAT|BARRAS} \qquad \illegible{1}\\
3 quai \unsure{Julliette|Juliette|Gailleton}}\\

\handwritten{Dr \unsure{MEUNIER|MECHAIN} 13 Av des 4 chemins \qquad \unsure{Médian|Médiau}\\
92330 SCEAUX}\\

\handwritten{\unsure{Dr} \unsure{Vivier|Viver} \qquad '' \qquad \illegible{1}}\\

\handwritten{\deleted{\unsure{Dr}} Roger \unsure{Vidal|Widal}. La Rochelle \qquad Roger \unsure{Widal|Vidal}.}\\

\handwritten{Dr \unsure{Vully|Vally}. 69002. \qquad \illegible{1}}\\
% archive sheet number: 131
